\section{Conclusion}
\label{sec:conclusion}

This thesis compared evaluated diffusion configurations and Qwen/Llama configurations on reasoning and planning tasks under shared downstream scoring conditions. The results do not support a universal ranking between the two model families. Instead, they show that the answer depends on the sampling budget and on the benchmark interface.

Countdown-cd4 gives the clearest example. The evaluated diffusion base configurations are strongest at estimated stochastic pass@1, but Qwen-Base and Llama-Base overtake them at high pass@k. Trip Planning shows a related but less clean pattern because parser compatibility and prompt format have a stronger effect. GSM8K is close to saturated at high $k$, so it is most useful here as a diagnostic for prompt format and answer extraction.

The evaluated diffusion configurations do not show the large-$k$ coverage advantage that motivated the study. Their strength appears mainly in deterministic or low-budget planning settings. The evaluated Qwen/Llama configurations, by contrast, often gain more from repeated stochastic sampling when many candidates are available. Per-question analyses show that the models solve overlapping but non-identical subsets of problems, which makes hybrid evaluation and selection methods a natural next step.

Future work should therefore focus less on declaring one generation interface superior and more on building practical systems that combine candidate generation with verification, routing, selection, or parser-aware repair.
